\chapter{Família IEC 61131 e o entorno}
\label{cap:familia-iec}

\begin{objetivos}
Ao final deste capítulo você deve ser capaz de:
\begin{itemize}
  \item identificar o assunto de cada parte da família IEC 61131;
  \item situar as quatro edições da parte 3 e o que mudou em cada uma;
  \item reconhecer as lacunas de cobertura do material antigo;
  \item explicar o papel da PLCopen, da IEC 61499 e do OPC UA;
  \item avaliar uma declaração de conformidade com a norma.
\end{itemize}
\end{objetivos}

\section{Como a norma está organizada}

A IEC 61131 é uma norma de controladores programáveis dividida em partes, preparadas e mantidas por
um grupo de trabalho específico (WG~7, sistemas de controle programáveis, do SC~65B do TC~65 da
IEC). Ela se chamava IEC 1131 antes da mudança de numeração da IEC.

\begin{table}[H]
  \centering
  \caption{Partes da família IEC 61131 e o assunto de cada uma.}
  \label{tab:partes-iec61131}
  \begin{tabular}{p{1.6cm} p{13.0cm}}
    \toprule
    Parte & Assunto \\
    \midrule
    1 & Informações gerais: definições de termos usados nas demais partes e as principais
        propriedades funcionais e características dos controladores \\
    2 & Requisitos de equipamento e testes: condições normais de serviço, requisitos construtivos,
        ensaios de tipo e compatibilidade eletromagnética \\
    3 & Linguagens de programação: arquitetura básica de software e linguagens do programa de
        controle \\
    4 & Orientações para o usuário: especificação, instalação, comissionamento e certificação \\
    5 & Comunicações: módulo de comunicação, blocos de função e mecanismos de conexão entre
        controladores e outros dispositivos \\
    6 & \textbf{Segurança funcional} \\
    7 & Programação de controle \emph{fuzzy} (FCL) \\
    8 & Orientações para aplicação e implementação das linguagens de programação \\
    9 & Interface de comunicação digital de ponto único para sensores e atuadores pequenos (SDCI,
        comercializada como IO-Link) \\
    10 & Formato de intercâmbio XML para exportação e importação de projetos IEC 61131-3 \\
    \bottomrule
  \end{tabular}
\end{table}

\begin{atencao}
Três partes acima merecem atenção de quem aprendeu pela apostila antiga. A \textbf{parte 6}
(segurança funcional) e as partes \textbf{9} e \textbf{10} não existiam quando o material de
referência foi escrito. Um texto que apresenta ``a família IEC 61131'' cobrindo apenas as partes 1 a
5, 7 e 8 está desatualizado, e a ausência da parte 6 é a mais grave: é ela que trata de controlador
em função de segurança.
\end{atencao}

\section{As edições da parte 3}

A parte 3 é a que interessa diretamente à programação, e ela evoluiu:

\begin{table}[H]
  \centering
  \caption{Edições da IEC 61131-3 e o que cada uma trouxe.}
  \label{tab:edicoes-61131-3}
  \begin{tabular}{p{2.4cm} p{12.2cm}}
    \toprule
    Edição & Marco \\
    \midrule
    1ª, 1993
      & Primeira publicação da padronização das linguagens de programação \\
    2ª, 2003
      & Edição usada como referência pelo material didático das apostilas de origem: define as cinco
        linguagens (LD, FBD, IL, ST e SFC) \\
    3ª, 2013
      & Descreve como implementar \textbf{programação orientada a objetos} no software aplicativo;
        \textbf{deprecia a Lista de Instruções (IL)} \\
    4ª, maio de 2025
      & Edição vigente: define três linguagens gráficas (LD, FBD, SFC) e uma textual (ST); a
        \textbf{Lista de Instruções foi removida} \\
    \bottomrule
  \end{tabular}
\end{table}

Duas consequências práticas para quem escreve material didático:

\begin{enumerate}
  \item dizer apenas ``conforme IEC 61131-3'' é insuficiente --- a edição precisa ser declarada;
  \item a IL deve aparecer como linguagem histórica e para leitura de código legado, nunca como
        recomendação de escrita (Capítulo~\ref{cap:linguagens}).
\end{enumerate}

\index{IEC 61131-3!edições}

\section{O entorno da norma}

A norma não vive sozinha. Três iniciativas completam o que ela não cobre:

\begin{figure}[H]
  \centering
  \includegraphics[width=0.95\linewidth]{12-ecossistema-norma.pdf}
  \caption{A norma e as iniciativas que a complementam.}
  \label{fig:ecossistema-norma}
\end{figure}

\paragraph{PLCopen.} Organização sem fins lucrativos que apoia e estende a norma. É dela a
biblioteca padronizada de blocos de função, o programa de certificação de ambientes de programação e
os esquemas XML de intercâmbio de projetos --- este último incorporado à família como parte 10 da
IEC 61131. Para o usuário, a PLCopen responde a uma pergunta prática: ``este programa roda no
controlador de outro fabricante se ambos forem conformes?'' --- e a resposta depende de quais
bibliotecas foram usadas.

\paragraph{IEC 61499.} Norma de blocos de função para \emph{controle distribuído}. Enquanto a IEC
61131-3 organiza o programa dentro de um controlador, a IEC 61499 organiza a aplicação espalhada por
vários dispositivos. É assunto de arquiteturas distribuídas, não substitui a 61131-3 no dia a dia da
máquina.

\paragraph{OPC UA.} Não é norma de programação, e sim camada de \emph{informação e comunicação}:
modelo de dados padronizado, independente de fabricante, com sessão segura. É o que permite que um
sistema supervisório ou um ERP leia dados de controladores de marcas diferentes sem escrever um
\emph{driver} para cada um. Detalhes no Capítulo~\ref{cap:iiot}.

\section{Como avaliar uma declaração de conformidade}

\begin{table}[H]
  \centering
  \caption{Perguntas para uma declaração de conformidade.}
  \label{tab:conformidade}
  \begin{tabular}{p{4.6cm} p{10.0cm}}
    \toprule
    Pergunta & Por que ela importa \\
    \midrule
    Qual edição da norma? & Há diferença relevante entre 2ª, 3ª e 4ª edições, inclusive sobre
      linguagens disponíveis \\
    Quais linguagens são suportadas? & ``Conforme'' não significa ``todas as linguagens'': o
      ambiente pode implementar apenas um subconjunto \\
    Há certificação de terceira parte? & Programa de certificação existe e é verificável; declaração
      do próprio fabricante não é o mesmo \\
    Os blocos usados são padronizados? & Blocos proprietários são o principal fator de dependência
      na migração \\
    Como é o endereçamento de E/S? & Endereçamento é justamente o que a norma não padroniza
      (Capítulo~\ref{cap:iec-versus-legado}) \\
    \bottomrule
  \end{tabular}
\end{table}

\begin{pratica}
Escolha um controlador e busque a declaração de conformidade à IEC 61131-3 disponível do fabricante.
Responda, com o documento em mãos:
(a) qual edição é citada?
(b) quais linguagens são efetivamente suportadas?
(c) existe certificação de terceira parte?
(d) quais blocos de função padronizados estão disponíveis?
Registre as respostas e as que ficaram sem resposta --- este último item é o mais informativo.
\end{pratica}

\begin{exercicios}
\begin{enumerate}[label=\alph*)]
  \item Qual parte da família trata de segurança funcional, e por que a ausência dela em um material
        de CLP é uma lacuna grave?
  \item Em que edição a Lista de Instruções foi depreciada e em qual foi removida?
  \item Explique a diferença de escopo entre IEC 61131-3 e IEC 61499.
  \item O que a PLCopen oferece que a norma não oferece diretamente?
  \item Um fornecedor apresenta um controlador ``conforme IEC 61131-3'' que implementa apenas Ladder.
        A afirmação é falsa? Justifique.
  \item Por que o formato de intercâmbio XML importa para quem precisa migrar de fabricante?
  \item Se uma planta tem controladores de três marcas diferentes, o OPC UA resolve o problema de
        portabilidade do \emph{programa}? Explique o que ele resolve e o que não resolve.
\end{enumerate}
\end{exercicios}
